\chapter*{Notation and Acronyms}
\noindent
The acronyms used in this report are listed below, followed by a note on the S7 addressing referred to in the design and evaluation chapters.

\vspace{0.5em}
\begin{center}
\small
\begin{tabular}{ll}
\hline
Acronym & Meaning \\
\hline
CARS & Criticality- and operation-Aware Response System (this project) \\
ICS & Industrial control system \\
OT & Operational technology \\
IT & Information technology \\
PLC & Programmable logic controller \\
HMI & Human-machine interface \\
EWS & Engineering workstation \\
SCADA & Supervisory control and data acquisition \\
SDN & Software-defined networking \\
OVS & Open vSwitch \\
DPI & Deep packet inspection \\
IDS & Intrusion detection system \\
FDI & False data injection \\
HIL & Hardware-in-the-loop \\
NAT & Network address translation \\
GUARD & The identity-binding, anti-spoof stage of the pipeline (Table~0) \\
MTTM & Time to mitigate (the reaction window from detection to enforcement) \\
REST & Representational state transfer (the controller's control interface) \\
TCP & Transmission control protocol \\
FC & Modbus function code \\
OB & Organisation block (PLC program unit) \\
DB & Data block (PLC memory area) \\
S7 & Siemens S7 communication protocol \\
PDU & Protocol data unit \\
TPKT & ISO transport packet header on TCP (the \texttt{03 00} bytes) carrying S7comm \\
COTP & Connection-oriented transport protocol (ISO 8073), the layer beneath S7comm \\
MBAP & Modbus application protocol header (carries the Modbus function code) \\
DPID & OpenFlow datapath identifier (per switch) \\
REAL & S7 IEEE~754 single-precision (32-bit) floating-point data type \\
TP, TN, FP, FN & True positive, true negative, false positive, false negative \\
\hline
\end{tabular}
\end{center}

\vspace{0.5em}
\noindent
S7 addressing follows the Siemens convention used throughout: \texttt{\%I} for process inputs and \texttt{\%Q} for process outputs at the bit level, \texttt{\%ID} and \texttt{\%QD} for the corresponding double-word (32-bit) analogue channels, which on this process carry REAL (IEEE~754 single-precision floating-point) values rather than double integers, so the level and valve values pass as native floats without \texttt{SCALE}/\texttt{NORM\_X} conversion, and \texttt{\%M} for internal memory bits (merkers). For example, \texttt{\%ID100} is the analogue level input from the tank sensor, \texttt{\%QD100} and \texttt{\%QD104} drive the fill and discharge valves, and \texttt{\%M0.1} and \texttt{\%M0.2} are the HMI start and stop bits.
